\begin{proof}[Proof of \cref{obj:lem:scale-sanity}]
Fix \(P\in\mathcal P_{d,\epsilon,M}^{\mathrm{unr}}\).  Since \(M\ge1\), we have \(M\ge0\).  For every supported cell \(k\),
\[
|\tau_k(P)|
=
|\mu_{1k}(P)-\mu_{0k}(P)|
\le |\mu_{1k}(P)|+|\mu_{0k}(P)|
\le M/2+M/2
=M .
\]
% lean: scale_sanity
For every cell \(k\), the product term in the all-cell formula of \cref{obj:def:ate-functional} satisfies
\[
|p_k(P)\tau_k(P)|\le p_k(P)M .
\]
Indeed, the cell-mass range gives either \(p_k(P)=0\) or \(p_k(P)>0\).  In the first case both sides are zero.  In the second case,
\[
|p_k(P)\tau_k(P)|
=p_k(P)|\tau_k(P)|
\le p_k(P)M,
\]
using the supported-cell bound above and \(p_k(P)>0\). % lean: scale_sanity
The finite cell masses satisfy
\[
\sum_{k\in[d]}p_k(P)\le1 .
\]
% lean: sum_cellMass_le_one
Hence, using the all-cell formula of \cref{obj:def:ate-functional},
\[
|\tau(P)|
=
\left|\sum_{k\in[d]}p_k(P)\tau_k(P)\right|
\le
\sum_{k\in[d]}|p_k(P)\tau_k(P)|
\le
\sum_{k\in[d]}p_k(P)M
=
\left(\sum_{k\in[d]}p_k(P)\right)M
\le M .
\]
% lean: scale_sanity
Therefore, for every supported cell \(k\),
\[
|\delta_k(P)|
=
|\tau_k(P)-\tau(P)|
\le |\tau_k(P)|+|\tau(P)|
\le M+M
=
2M .
\]
% lean: scale_sanity

The radius-two membership follows by keeping the same law and the same consistency, exchangeability, overlap, mean-normalization, and second-moment witnesses, using the numerical facts \(0\le2\le2\), and adding the just-proved supported-cell bound
\[
|\delta_k(P)|\le 2M=2\cdot M
\qquad (p_k(P)>0)
\]
as the approximate-homogeneity witness. % lean: scale_sanity
Thus every law in \(\mathcal P_{d,\epsilon,M}^{\mathrm{unr}}\) is represented in \(\mathcal P_{d,\epsilon,M,2}\).  Conversely, a member of \(\mathcal P_{d,\epsilon,M,2}\) carries all fields required for \(\mathcal P_{d,\epsilon,M}^{\mathrm{unr}}\); forgetting the homogeneity and radius fields leaves the same law. % lean: scale_sanity
\end{proof}
